\chapter{Úvod}
Testovanie je jedna z najdôležitejších fáz vývoja softvéru, pretože dôsledky nekvalitného softvéru sú nákladné. Nekvalitný softvér spôsobil v roku 2022 v USA náklady odhadované na skoro dva a pol bilióna amerických dolárov[tu pridaj citacu krasner 2022]. Jedným z prístupov ku testovaniu je Data-Driven testing (DDT), testovanie založené na princípe oddelenia testovacích dát od testovacej logiky.

Efektívne testovanie stojí na kvalite testovacích dát, hlavne preto, že pri nedostatočne domyslených dátach môže dôjsť k neodhaleným chybám, ale aj k falošnému pocitu dobre otestovaného produktu, čo môže spôsobiť stratu časových a finančných zdrojov.
Vytváranie testovacích dát je časovo náročné a vyžaduje znalosť schémy, ale aj charakteru reálnych dát.

Pre uľahčenie práce existujú nástroje na generovanie náhodných testovacích dát. Jedným z nich je Dataster, nástroj, ktorý generuje štruktúrované testovacie dáta pre relačné databázy na základe konfiguračného receptu. Ručné vytvorenie receptu vo forme YAML vyžaduje znalosť dokumentácie programu Dataster, ako aj znalosť schémy databázy, vzťahov medzi tabuľkami a charakteru uložených hodnôt. Náročnosť tohto procesu rastie so zložitosťou databázy. Pre používateľa s nízkou znalosťou konkrétnej databázy predstavuje tvorba receptu netriviálnu manuálnu prácu.

Potenciálnym riešením je automatizácia tejto úlohy s pomocou veľkých jazykových modelov (LLM), ktoré majú schopnosť interpretovať štruktúrované dáta a transformovať ich do iných formátov. Táto práca využíva lokálne LLM, keďže vstupom sú reálne dáta z databázy a lokálny beh zabezpečuje nielen súkromie dát, ale aj nízke náklady.

Cieľom tejto práce je navrhnúť a implementovať nástroj, ktorý sa pripojí k databáze, extrahuje jej schému a vzorku uložených dát a pomocou lokálnych LLM automaticky vygeneruje recept pre nástroj Dataster. Ďalším cieľom je zhodnotiť kvalitu generovaných receptov naprieč rôznymi modelmi a konfiguráciami. Kvalita sa meria od štrukturálnej správnosti vygenerovaného receptu (YAML) cez akceptovanie nástrojom Dataster až po vykonateľnosť výsledného SQL kódu. Vyhodnotenie prebieha na sade testovacích databáz rôznej zložitosti.

Kapitola 2 približuje definíciu testovania a popis data-driven testing. Kapitola 3 definuje LLM, ich využitie pri generovaní dát, predstavuje jednotlivé LLM modely a skúmanie alternatív. Kapitola 4 približuje nástroj Dataster, jeho využitie, limitácie, popis konfiguračného receptu a YAML súborov. Kapitola 5 predstavuje vysokoúrovňový návrh programu. Kapitola 6 sa zaoberá koncovou implementáciou. Kapitola 7 vysvetľuje proces experimentov, sadu testovacích databáz, úrovne validácie vygenerovaných receptov a porovnávané modely a konfigurácie. Kapitola 8 prezentuje a porovnáva výsledky experimentov. Kapitola 9 objasňuje limitácie programu a možné rozšírenia v budúcej práci. Na záver kapitola 10 zhŕňa dosiahnuté výsledky a porovnáva ich s cieľmi práce.



\chapter{Testovanie softvéru}

Cieľom tejto kapitoly je vymedziť pojem testovanie, jeho význam a dôsledky jeho zanedbania. Následne kapitola definuje termín Data-Driven testing (DDT), spôsob testovania, ktorý oddeľuje testovacie dáta od testovacej logiky a tým kladie zvýšené nároky na kvalitu a prípravu testovacích dát. Práve náročnosť prípravy dát je jeden z hlavných východiskom tejto práce.

\section{Význam testovania}
Testovanie je možné chápať ako proces vyhodnocovania softvéru na základe pozorovania jeho behu [AmmannOffutt2016]. Jeho cieľom je odhaliť chyby a overiť, či softvér spĺňa špecifikované požiadavky a potreby používateľov. [Myers 2011].

\section{Dôsledky nedôkladného testovania}
Každý softvér je v nejakej fáze vývoja chybný. Tieto chyby môžu mať za následok závažné straty, začínajúc od finančných strát až po straty na ľudských životoch. Existuje množstvo zaznamenaných prípadov kde nedôkladné testovanie viedlo až ku katastrofickým následkom. Jeden z týchto prípadov je explodovanie rakety  Ariane 5 (let 501), 37 sekúnd po jej štarte v roku 1996. Dôvodom bola neošetrená výnimka pri konverzii 64-bitového čísla s plávajúcou desatinnou čiarkou (float) na 16-bitové celé číslo so znamienkom (signed integer) v navigačnom systéme keďže prevzatý modul nebol testovaný pre letové parametre novej rakety. Toto viedlo k strate odhadovanej na približne 370 miliónov dolárov a ročnému zdržaniu programu Ariane 5. [LeLann1996]

Ďalším príkladom je Therac-25, prístroj určený na rádioterapiu, ktorý bol používaný v USA a Kanade od roku 1983 až do ich stiahnutia v roku 1987. Chybný softvér spôsobil šesť prípadov masívneho predávkovania radiáciou, z toho tri úmrtia. Podobné chyby boli nájdené aj v predošlom modeli Therac-20, ten mal ale hardvérové poistky, ktoré pri rovnakej softvérovej chybe zabránili ožiareniu. Therac-25 ich nahradil čisto softvérovou kontrolou. [Leveson1993].

Podobných prípadov je ale viac. Poukazujú na dôležitosť testovania, nevyhnuteľné pre bezpečnosť softvéru.  Okrem katastrofických zlyhaní generuje nekvalitný softvér spôsobuje aj priebežné ekonomické náklady. [Krasner2022]


\section{Úrovne testovania}
Testovanie prebieha na rôznych úrovniach, pričom každá overuje iné aspekty systému. [Roman2018]

\begin{itemize}
    \item \textbf{Jednotkové testovanie (unit testing)} - testovanie najmenších jednotiek kódu ako funkcie a metódy.
    \item \textbf{Testovanie modulov (module testing)} - testovanie modulu (kolekcie jednotiek) v izolácii. Hľadajú chyby v komunikácii jednotiek v rámci modulu. 
    \item \textbf{Integračné testovanie (integration testing)} - overenie rozhrania a komunikácie medzi modulmi. 
    \item \textbf{Testovanie systému (system testing)} - kontroluje softvér ako celok. Hľadajú sa chyby v návrhu a špecifikácii.
    \item \textbf{Akceptačné testovanie (acceptance testing)} - overenie, či softvér spĺňa požiadavky zákazníka a či spĺňa svoj primárny účel.
\end{itemize}

Táto práca sa dotýka predovšetkým vyšších úrovní, keďže generované testovacie dáta slúžia na plnenie databáz pri integračnom a systémovom testovaní aplikácií nad nimi.

\section{Automatizované testovanie}
Pri väčších projektoch sa manuálne testovanie stáva neefektívnym, nielen z časového, ale aj ekonomického hľadiska. Používa sa preto automatizované testovanie, pri ktorom testy vykonávajú špecializované nástroje a frameworky. Tento spôsob testovania sa stáva nevyhnutnou súčasťou vývoja. Automatizované testy sa spúšťajú opakovane a s rôznymi vstupmi, čím sa príprava testovacích dát stáva samostatnou, opakovanou potrebou. Jedným z prístupov testovania, ktorý oddeľuje testovacích dát od testovacej logiky, je Data-Driven testing (DDT), na ktorý sa zameriava táto práca.

\section{Data-Driven testing}
Základnou myšlienkou je, že rovnaký testovací skript možno spustiť s rôznymi sadami vstupných dát, čím sa zvyšuje pokrytie testov bez nutnosti duplikovania kódu. [Laukkanen]

\subsection{Hlavné charakteristiky} 
Oddelenie dát od logiky - testovacie dáta sú uložené v externých súboroch ako napríklad Excel, CSV, YAML, nie v kóde testu.
\begin{itemize}
    \item \textbf{Znovupoužiteľnosť} - rovnaké testy možno spustiť s rôznymi sadami dát.
    \item \textbf{Čitateľnosť} - testovacie dáta musia byť ľahko čitateľné aj osobou bez znalostí programovania [Laukkanen]
    \item \textbf{Jednoduchá údržba} - možné pridanie nového testovacieho prípadu bez zmeny kódu.
\end{itemize}

Data-Driven testing prináša mnoho výhod, ako zvýšenie pokrytia testov a jednoduchá údržba testov, má však aj svoje nevýhody. Je potrebné testovacie dáta neustále udržiavať a príprava kvalitných testovacích dát môže byť náročná. Kvalita testovania je limitovaná kvalitou testovacích dát a to je motivácia pre využitie automatizovaného generovania testovacích dát.


\chapter{LLM}

\chapter{DATASTER}


\chapter{Návrh}

Táto kapitola sa zaoberá vysokoúrovňovým návrhom nástroja na automatizované generovanie receptov pre nástroj Dataster. Následne približuje celkovú architektúru, komponenty návrhu a tok dát medzi nimi. LLM tvorí jadro návrhového riešenia a preto je pozornosť venovaná spôsobu využitia veľkých jazykových modelov. Kapitolu uzatvára princíp princíp validácie a opravy výstupov. Konkrétna implementácia návrhu je obsiahnutá v kapitole 6.

\section{Architektúra}
Vstupom nástroja sú identifikačné údaje pre pripojenie k relačnej databáze. Na začiatku sa program pripojí k databáze, extrahuje tabuľky, informácie o nich, teda stĺpce a ich dátové typy. Po získaní týchto informácií sa z tabuľky vytiahnu rôzne vzorky, ktoré zahŕňajú krajné aj typické hodnoty každého stĺpca, aby mal jazykový model k dispozícii reálne hodnoty z tabuľky.

V ďalšej fáze LLM ako prvé vygeneruje types sekciu, keďže na tejto sekcii závisia ostatné sekcie YAML. Následne sa vygenerujú ďalšie sekcie. Sekcie schema, generate a constraints sa generujú nezávisle, každá samostatným volaním modelu s types sekciou ako vstupom. 

Sekcia types prechádza opravami ešte pred generovaním ostatných sekcií, keďže slúži ako ich vstup a chyby v nej by sa šírili ďalej. Zo všetkých vygenerovaných sekcií  sa poskladá jeden predbežný YAML, ktorý je validovaný a v prípade chýb opravený. Na záver je recept opäť skontrolovaný a uložený.

\begin{figure}[h]
    \centering
    \includegraphics[width=1\linewidth]{obrazky-figures/diagram1.pdf}
    \caption{tu nieco pridaj}
    \label{fig:diagram1}
\end{figure}

\section{Fázy generovania}
Samotné generovanie sekcií je rozdelené do niekoľkých fáz, hlavne z dôvodu izolácie chýb. Pri zlyhaní nie sú ostatné sekcie stratené a chyba je priraditeľná fáze, v ktorej vznikla. Malé modely zvládajú vymedzené úlohy s few-shot vzormi lepšie ako komplexné viacčasťové výstupy, kde s rastúcou zložitosťou požadovaného výstupu rastie výskyt formátových chýb [Pysmennyi2025]. Rozdelenie úloh šetrí pamäť. Každý krok spracováva len nutné informácie, čo je pri 8GB VRAM zásadná výhoda.

Pre fázu types je vstupom sterilizovaná schéma a vzorky dát, ktorý obsahuje informácie o primárnom kľúči tabuľky, dátových typoch stĺpcov ako aj niekoľko odlišných vzoriek. Ide o interpretačne najnáročnejšiu fázu a jej kvalita vygenerovaného výstupu priamo ovplyvňuje ostatné fázy, preto sa jej návrhu venuje najväčšia pozornosť. 

Fáza schema prevezme vygenerovanú sekciu types. Jej úlohou je mapovanie stĺpcov a vzťahov cudzích kľúčov. Hoci je mapovanie z veľkej časti mechanické, rozpoznanie vzťahov medzi tabuľkami vyžaduje interpretáciu názvov a hodnôt, preto je úloha zverená modelu.

Pri fáze generate sa určujú jednotlivé počty generovaných riadkov pre každú tabuľku. Logický počet sa odvíja od rozsahu hodnôt primárneho kľúča.

Ako posledná fáza je constraints, ktorá podľa vstupu types odvodzuje logické obmedzenia hodnôt, ktoré nie sú zachytené v schéme databázy, ale vyplývajú z kontextu hodnôt, ako napríklad rozsahy veku alebo nezáporné ceny. Tieto obmedzenia nie je možné mechanicky vyčítať, model ich musí vyvodiť zo vzoriek.

Toto poradie je dôležité pretože types je spoločným vstupom ostatných troch generovaných sekcií. Types definuje množinu tabuliek, stĺpcov a ich typov, na ktorú sa ostatné sekcie odkazujú. Ostatné fázy sú od seba  navzájom nezávislé a preto sa volajú oddelene, chyba v jednej z týchto častí neruší ostatné. 

\section{Využitie veľkých jazykových modelov (LLM)}
Jazykový model v návrhu nefunguje ako autonómny agent, ktorý sám rozhoduje o ďalšom kroku, poradie fáz vrátane validácie a opráv sú pevne riadené programom. AI agenti sú síce populárni čo sa týka zabezpečovania kvality softvéru [Pysmennyi2025], ale pre tento program to bolo vyhodnotené ako nevhodné hneď z niekoľkých dôvodov. Po prvé, experimentálne vyhodnotenie vyžaduje reprodukovateľnosť, rovnaký vstup musí prejsť tou istou cestou. S agentom pribudne ďalšia neistota, pretože pri každom spustení sa môže líšiť nielen odpoveď modelu, ale aj počet a poradie krokov. Po druhé, modely s veľkosťou okolo 7 miliárd parametrov, ktorú si vynucuje dostupná pamäť GPU (8GB VRAM) nemajú kapacitu na spoľahlivé plánovanie [Valmeekam2022]. Po tretie, keď má každá fáza svoje pevné miesto, zlyhanie je priraditeľné k fáze v ktorej sa objavilo, pri slučke tvorenej agentom sa chyba stráca v trajektórii.

Modely tak rozhodujú výhradne o obsahu generovaných sekcií, nie o štruktúre procesu. Všetko, čo je možné vykonať algoritmicky, vykonáva kód.

\section{Validácia a kontrola}
Priebeh validácie znázorňuje obrázok (pridaj číslo obrázka), teda pri validácii môžu nastať 3 prípady. Bud je vygenerovaný YAML bez chýb a pokračuje ku ďalšej fáze, alebo prechádza cyklom opráv. V prípade, že kód narazí na chybu, ktorá pretrváva aj po N cykloch opráv, kód končí s chybovým hlásením.

\begin{figure}[h]
    \centering
    \includegraphics[width=0.5\linewidth]{obrazky-figures/diagram2.pdf}
    \caption{tu nieco pridaj}
    \label{fig:diagram2}
\end{figure}


Výstupy jazykových modelov sú formátovo nestabilné a menšie modely majú zdokumentované problémy s dodržaním konzistentného výstupného formátu [Pysmennyi2025]. Preto bola validácia a oprava výstupov navrhnutá ako integrálna súčasť procesu, nie ako dodatočná kontrola.

\section{Rozdelenie práce}
Využitie LLM pre každú časť procesu, od extrakcie schémy cez generovanie až po opravy by bolo nielen zbytočné, ale aj pomalé. Návrh túto prácu delí medzi LLM a samostatný kód podľa povahy úlohy. 

LLM sa stará o generáciu YAML častí ako types, schema, generate a constraints, ale z generovania vynecháva sekciu include a SQL požiadavky. LLM dostáva úlohy pri ktorých je dôležitá interpretácia a neexistuje presný algoritmus.

Samotný kód dostáva úlohy na mechanickej báze, kde je nutný presný postup a odchýlka od neho je chybou. Stará sa o pripojenie k databáze, extrakciu schémy a vzoriek dát ako aj poskladanie týchto informácií do jedného balíka, ktorý je poslaný do LLM. Ako posledné si prevezme vygenerovaný výsledok, ten skontroluje, opraví a uloží. 

V počiatočných fázach návrhu bol testovaný aj variant fázy, ktorý sám opravoval chybný recept, lenže opravy boli pomalé a nespoľahlivé a nakoniec sa oprava kódom ukázala ako lepšia voľba. 

Pôvodný návrh zo zadania počítal s LLM aj pri generovaní SQL dopytov ako aj generovaní YAML include sekcie. Obe úlohy sa ale ukázali byť plne programovateľné a keďže LLM nepridávalo hodnotu ani  jednej zo spomenutých častí a vytváral zbytočné chyby, implementácia týchto častí bola zverená kódu.

 
